\section{Explicit normal forms at singular strata}
\label{sec:singular-normal-forms}
The semialgebraic envelope does not by itself compute a singularity.
We give two complementary computations. The first identifies the
leading constant at real-rooted multiplicities by a finite quartic
cone in the recovered coefficient algebra. The second treats an
intersection of channel collapse, normalization failure, component
equality, and internal multiplicity directly from the model equations.
At the latter intersection the modulus has a nonzero constant term;
its transverse regular condition has an explicit two-regime law.

\subsection{A profiled quartic cone for real-rooted multiplicities}
Assume the hypotheses of Proposition~\ref{prop:coefficient-recovery},
with all roots interior, and at least one internal multiple root.
For each component let $r_g$ run over its distinct root values, and
let $m_g$ be their multiplicities. A group $g$ includes its component
index, even when another component has the same scalar root.
Let $\mathcal M=\{g:m_g\ge2\}$. In the ambient coefficient tangent
space, define $D_g$, $g\in\mathcal M$, to have zero channel and weight
coordinates and component-polynomial variation
\[
 (D_g)_b(z)=f_b(z)/(z-r_g)^2
\]
in its own component, and zero variations in other components.
Let $\mathcal T_0$ be the linear space generated by all channel and
weight tangents and by the vectors with polynomial variation
$f_b(z)/(z-r_g)$, one for every distinct root group. These last
vectors are the cluster-centroid tangents, including the simple roots.

Put $\langle u,v\rangle_I=\langle J_cu,J_cv\rangle_{H_{P,w}}$.
This is positive definite by Proposition~\ref{prop:coefficient-recovery}.
Let $\Pi_0$ be its orthogonal projection onto $\mathcal T_0$, and set
\begin{equation}\label{eq:profiled-split-information}
 S_{gh}=\langle (I-\Pi_0)D_g,(I-\Pi_0)D_h\rangle_I,
                         \qquad g,h\in\mathcal M.
\end{equation}
Here the identity in $I-\Pi_0$ is an operator, not the metric matrix.
The matrix $S$ is a Schur complement of the recovered coefficient
Fisher information. In particular it profiles out the stochastic
factors and the cluster centroids, rather than holding nuisance
parameters fixed.

For a vector $x$ with one coordinate per root, set
$v_g(x)=\tfrac12\sum_{i\in g}x_i^2$. Define the compact set
\begin{equation}\label{eq:quartic-cone-section}
 \mathcal X_S=\left\{x:
 \begin{array}{l}
 x_i=0\ \text{on simple groups},\quad
 \sum_{i\in g}x_i=0\quad(g\in\mathcal M),\\
 v(x)^{\mathsf T}S v(x)\le4
 \end{array}\right\}.
\end{equation}
As before, $\Gamma$ permutes all root coordinates with the same base
scalar value, including coordinates belonging to different components.

\begin{theorem}[The real-rooted singular leading constant]
\label{thm:quartic-modulus}
The matrix $S$ in \eqref{eq:profiled-split-information} is positive
definite, and
\begin{equation}\label{eq:quartic-modulus}
 \omega_P(t)=C_S\sqrt t+o(\sqrt t),\qquad
 C_S=\max_{x,y\in\mathcal X_S}d_\Gamma(x,y)>0.
\end{equation}
All quantities defining $\mathcal X_S$ are recovered from $P$;
coefficient-coordinate changes leave the construction invariant, and
component relabelling only permutes its groups. Thus the exponent and
leading constant are determined by a finite profiled quartic problem,
not by elimination over the original closed parameter space.
If there is exactly one multiple group, it has multiplicity two, and
its base value is not shared by any other component, then, writing
$S=s>0$, the constant is
\begin{equation}\label{eq:double-root-constant}
 C_S=\sqrt2\,s^{-1/4}.
\end{equation}
\end{theorem}
\begin{proof}
For each component the rational functions
$1/(z-r_g)$ and $1/(z-r_g)^2$ have linearly independent principal
parts at their distinct poles. Multiplication by $f_b$ shows that the
$D_g$ are independent modulo $\mathcal T_0$. Injectivity of $J_c$
then makes $S$ positive definite. It follows that
$v^{\mathsf T}Sv\le4$ bounds every $v_g$ and every root coordinate;
hence \eqref{eq:quartic-cone-section} is compact.

We prove the limiting set assertion underlying
\eqref{eq:quartic-modulus}. Proposition~\ref{prop:coefficient-recovery}
gives coefficient displacement $O(t)$ for every member of
$\mathcal B_P(t)$. Analytic factorization into the pairwise coprime
cluster factors, as in Lemma~\ref{lem:anchored-real-roots}, therefore
gives $O(t)$ displacement of each cluster coefficient. Write the real
root displacements in group $g$ as $y_{gi}$. The first two elementary
symmetric coefficients give
\[
 \sum_i y_{gi}=O(t),\qquad
 \sum_i y_{gi}^2=O(t).
\]
Simple roots have displacement $O(t)$. After a subsequence,
$x_{gi}=y_{gi}/\sqrt t$ converges, its sum is zero, and the first
coefficient divided by $t$ also converges. Elementary symmetric
terms of degree at least three are $O(t^{3/2})=o(t)$. Consequently
every limiting ambient coefficient displacement divided by $t$ has
the form
\begin{equation}\label{eq:split-normal-form}
 a_0-\sum_{g\in\mathcal M}v_g(x)D_g,
                               \qquad a_0\in\mathcal T_0.
\end{equation}
The Hellinger expansion in the proof of
Theorem~\ref{thm:regular-modulus-bridge} is valid on this ambient
coefficient chart. Profiling \eqref{eq:split-normal-form} over $a_0$
therefore gives $v(x)^{\mathsf T}Sv(x)\le4$. This proves the outer
limiting-set inclusion for the rescaled target.

Conversely take $x\in\mathcal X_S$ with strict inequality. Choose
$a_0\in\mathcal T_0$ attaining the orthogonal projection minimum.
Realize its channel and weight coordinates by linear perturbations
of size $t$. In group $g$ place the roots at
\[
 r_g+\sqrt t\,x_{gi}+t b_g,
\]
where the common centroid shift $b_g$ is selected to give the
corresponding $f_b/(z-r_g)$ coefficient of $a_0$; use the analogous
linear shift for each simple root. The zero-sum condition gives
exactly \eqref{eq:split-normal-form} at order $t$. Higher terms are
$o(t)$. These are real admissible interior roots, and the strict
stochastic, rank, normalization, and tuple conditions persist.
Their Hellinger distance divided by $t$ tends to
$\tfrac12\sqrt{v(x)^{\mathsf T}Sv(x)}<1$. Points on the boundary
are approached by replacing $x$ by $(1-\varepsilon)x$ and then
letting $\varepsilon\downarrow0$. Thus the target sets, divided by
$\sqrt t$ and taken modulo $\Gamma$, converge in Hausdorff distance
to $\mathcal X_S/\Gamma$. Matching between distinct base values is
excluded by their fixed gaps. Taking diameters proves
\eqref{eq:quartic-modulus}.

A symmetric split in any multiple group gives a nonzero member of
$\mathcal X_S$, while $0$ also belongs to it, so $C_S>0$. The
coordinate invariance follows from the pullback metric and the
intrinsic principal-part directions; permutations just relabel them.
In the final case $x=(-a,a)$ in the double group and zero elsewhere,
with $a\ge0$ after quotienting. Its constraint is $s a^4\le4$ and
its quotient distance is $|a-a'|$. The diameter of this interval is
$\sqrt2\,s^{-1/4}$, proving \eqref{eq:double-root-constant}.
\end{proof}

\begin{theorem}[Uniform crossover at a splitting double root]
\label{thm:double-root-crossover}
Suppose just one component has a double root $r$, all other component
roots are simple and separated from $r$, and the strict hypotheses of
Proposition~\ref{prop:coefficient-recovery} hold. Replace this double
factor by $(z-r)^2-\Delta^2$, leaving all other parameters fixed,
and denote the datum by $P_\Delta$. For sufficiently small fixed
$\Delta_0,t_0>0$, uniformly for $0\le\Delta\le\Delta_0$ and
$0<t\le t_0$,
\begin{equation}\label{eq:double-root-crossover}
 \omega_{P_\Delta}(t)\asymp
             t+\frac{t}{\Delta+\sqrt t}.
\end{equation}
For $\Delta>0$ its regular Fisher condition is of order
$1+\Delta^{-1}$. The linear regime $t\ll\Delta^2$ and the singular
regime $t\gg\Delta^2$ therefore join at $t\asymp\Delta^2$.
At $\Delta=0$ the exact leading constant is
\eqref{eq:double-root-constant}.
\end{theorem}
\begin{proof}
The coefficient recovery map and its inverse differential are bounded
uniformly on a small compact neighbourhood of $P_0$. Coprime cluster
factorization is uniform there as well. Write the exceptional monic
quadratic as $(z-r-\mu)^2-v$, with $v\ge0$. In coefficient coordinates
this is an analytic change of variables even at $v=0$.
A Hellinger ball of radius $t$ about $P_\Delta$ consequently has
$|\mu|\le Ct$, $|v-\Delta^2|\le Ct$, and $O(t)$ changes of all the
other roots. Uniform exclusion of distant representatives follows
from the recovery proposition and compactness. For two such quadratic
factors, their sorted-root distance is at most the difference of
centres plus $|\sqrt v-\sqrt{v'}|$. The latter is bounded by
$C t/(\Delta+\sqrt t)$, by considering separately
$t\le c\Delta^2$ and its complement. This proves the upper bound.

For the lower bound increase $v$ from $\Delta^2$ to
$\Delta^2+ct$, with all other coordinates fixed. For sufficiently
small fixed $c>0$ this alternative has Hellinger distance at most $t$,
and its target displacement is
\[
 \sqrt{\Delta^2+ct}-\Delta
       =\frac{ct}{\sqrt{\Delta^2+ct}+\Delta}
       \asymp\frac{t}{\Delta+\sqrt t}.
\]
A common shift of the quadratic roots by $ct$ supplies the linear
term. The larger of the two lower bounds controls their sum up to a
factor two. Interior margins preserve admissibility. For $\Delta>0$
root differentiation in $(\mu,v)$ gives derivatives $1$ and
$\pm(2\Delta)^{-1}$, whereas the coefficient Fisher form is uniformly
positive and bounded. This proves the condition estimate and all the
stated transitions.
\end{proof}

\subsection{A binary intersection of normalization and channel defects}
Here $k=m_1=m_2=2$, the weights are $1/2$, and $\ell\ge2d+1$.
For $0\le a,b<1$, put
\[
 U_a=\tfrac12\begin{pmatrix}1+a&1-a\\1-a&1+a\end{pmatrix},\qquad
 V_b=\tfrac12\begin{pmatrix}1+b&1-b\\1-b&1+b\end{pmatrix},
 \qquad f_\pm(z)=(z-h\mp x)^d,
\]
with $h\pm x$ interior. Define the observable binary moments
$u=P_{11}+P_{12}-P_{21}-P_{22}$,
$v=P_{11}+P_{21}-P_{12}-P_{22}$, and
$w=P_{11}+P_{22}-P_{12}-P_{21}$. Direct substitution gives the exact
normal form
\begin{equation}\label{eq:binary-normal-form}
 (u_j,v_j,w_j)=(a\chi_j,b\chi_j,ab),\qquad
 \chi_j=\frac{(T_j-h-x)^d-(T_j-h+x)^d}
              {(T_j-h-x)^d+(T_j-h+x)^d}.
\end{equation}
Let $P_*$ be the constant matrix with every entry $1/4$.

\begin{theorem}[Collapse normal form and transverse phase law]
\label{thm:binary-collapse}
Near $a=b=x=0$, the observation ideal in analytic coordinates is
\begin{equation}\label{eq:binary-defect-ideal}
                  (ab,ax,bx).
\end{equation}
In particular, uniformly for $h$ in an interior compact interval,
\[
 h_w(P_{a,b,x,h},P_*)\asymp |ab|+|ax|+|bx|.
\]
The intersection $a=b=x=0$ has two rank-one channels, equal complete
component polynomials, $\tau(P_*)=0$, and internal multiplicity $d$
when $d\ge2$. Its whole-model modulus is exactly
\begin{equation}\label{eq:collapse-modulus}
                     \omega_{P_*}(t)=D-L\qquad(t\ge0).
\end{equation}
For arcs $a=A s^p+o(s^p)$, $b=B s^q+o(s^q)$,
$x=X s^r+o(s^r)$ with $A,B,X>0$, $p,q,r>0$, and $h=h_0$,
the observation order is
$\min\{p+q,p+r,q+r\}$, whereas the target displacement from the
selected repeated-root representative is of order $s^r$.

In degree one, at every $a,b,x>0$ in a sufficiently small fixed
parameter box, the regular condition for the \emph{entire} binary
closed model, not merely for the displayed subfamily, satisfies
\begin{equation}\label{eq:binary-condition}
 \mathfrak c_{H_{P,w}}(P)\asymp\eta_1(P)^{-1}
       \asymp \frac1{ab}+\frac1{x\sqrt{a^2+b^2}}.
\end{equation}
The constants are uniform in $a,b,x$. Along the preceding power arcs
its divergence exponent is
\begin{equation}\label{eq:binary-phase}
                 \max\{p+q,\ r+\min(p,q)\}.
\end{equation}
The phase boundary is $r=\max(p,q)$. At each regular datum the leading
observation-ball constant is four times this Fisher condition,
by Theorem~\ref{thm:regular-modulus-bridge}.
\end{theorem}
\begin{proof}
The binary moment transformation is invertible, with
$P_{ij}=\tfrac14(1+\varepsilon_i u+\varepsilon_j v+
\varepsilon_i\varepsilon_j w)$, where $\varepsilon_1=1$ and
$\varepsilon_2=-1$. Equation~\eqref{eq:binary-normal-form} is obtained
by averaging the two rank-one component matrices. At each clock,
\[
 \chi_j=-\frac{d x}{T_j-h}+O(x^3).
\]
Thus $\chi_j/x$ is an analytic unit near the intersection. The ideal
generated by all moment differences from $P_*$ is exactly
\eqref{eq:binary-defect-ideal}. Positive cell floors make Hellinger
and moment norms uniformly equivalent. This proves the stated
norm and arc orders without quantifier elimination. More precisely,
if $s_0=\min\{p+q,p+r,q+r\}$ and
$C_h=\sum_jw_j(T_j-h_0)^{-2}$, the squared leading Hellinger coefficient
is
\[
 \tfrac14\left[
 \mathbf1_{p+q=s_0}(AB)^2+
 d^2C_h\{\mathbf1_{p+r=s_0}(AX)^2+
                   \mathbf1_{q+r=s_0}(BX)^2\}\right].
\]
The three moment directions are orthogonal at $P_*$, so no leading
cancellation is suppressed in this expression.

All coordinates of any two spectral targets lie in $[L,D]$, giving
$\omega_{P_*}(t)\le D-L$. With both channels equal to their collapsed
columns, taking every component polynomial to be $(z-L)^d$ or every
one to be $(z-D)^d$ gives the same datum $P_*$. These two targets
attain $D-L$, proving \eqref{eq:collapse-modulus}. More generally any
constant datum admitting the stochastic factorization of the model
has this same spectral diameter, by keeping its factors and weights
and choosing the common polynomial at either endpoint. Thus a positive
vanishing exponent is not the appropriate invariant at this
intersection: the computed constant exact-fibre term persists.

It remains to establish the whole-model transverse assertion. For
$d=1$, $P(T)=A+B/(T-h)$, with $A$ having binary moments $(0,0,ab)$
and $B$ having zero total and moments $(-ax,-bx,0)$. The affine
normalization calculation gives
$\tau(P)\asymp\|B\|_F\asymp x\sqrt{a^2+b^2}$.
The rank-one residue formula gives
\[
 B_1(P)=\sum_{i=1}^2
       2\|V_b^{-\mathsf T}e_i\|_2\|e_i^{\mathsf T}U_a^{-1}\|_2
                       \asymp (ab)^{-1}.
\]
Cell floors and clock factors are uniform, so the global observable
bound and Theorem~\ref{thm:exact-condition} give the upper estimate
in \eqref{eq:binary-condition} for all model tangents.

Varying $h$ with $a,b,x$ fixed moves both roots at unit speed and
has observation derivative of Fisher norm comparable to
$x\sqrt{a^2+b^2}$. This gives the normalization lower bound. For the
other lower bound keep $q=z-h$ and $A$ fixed and add to $B$ a matrix
with zero total and binary moments $(0,0,\varepsilon)$. The determinant
of the perturbed numerator, writing $y=z-h$, is
\[
                 \tfrac14\{ab\,y^2+\varepsilon y-ab\,x^2\}.
\]
Both roots have derivative $-1/(2ab)$ at $\varepsilon=0$, while the
observation derivative has uniformly bounded Fisher norm. This is an
admissible model tangent: the pencil has two distinct real roots at
the base; its analytic rank-one spectral factors are strictly positive
there and remain so for sufficiently small $\varepsilon$. Their
positive total, row, and column sums give strict weights and stochastic
columns, and the total numerator remains $q$. The two roots remain
interior. The size of this admissible interval may depend on the base
point, which is sufficient for the differential assertion. The larger
of the two lower bounds controls their sum. This proves
\eqref{eq:binary-condition}; substituting the power arcs proves
\eqref{eq:binary-phase}. The complete polynomials have distinct roots
for $x>0$, so the collision group is trivial and the final assertion
follows from \eqref{eq:regular-modulus-constant}.
\end{proof}

For $a,b,x>0$ in the degree-one normal form, these transverse defect
coordinates are themselves recoverable from observations. The three
clock ratios give $h$; the moment residues give $ax$ and $bx$, while
the constant joint moment gives $ab$. Thus
$x^2=(ax)(bx)/(ab)$ and similarly for $a^2,b^2$. Their positive square
roots fix the displayed branch. The exponents in
\eqref{eq:binary-phase} therefore describe observable relative orders,
not inaccessible choices of latent normalization. At the intersection,
\eqref{eq:collapse-modulus} replaces the diverging differential law.
The variational envelope remains valid on other singular faces; the
explicit computations here do not identify defect dimensions alone
with a universal classification of all of them.
